%!TEX root = tfg.tex
\chapter{Conclusiones}

\begin{quotation}[Poet]{Dante Alighieri (1265--1321)}
Consider your origin: you were not made to live as brutes, but to follow virtue and knowledge.
\end{quotation}

\begin{abstract}
Este capítulo recoge la valoración final del proyecto, sintetizando el aprendizaje
obtenido durante el desarrollo y señalando las líneas de mejora que quedan abiertas.
\end{abstract}

\section{Informe post-mortem}
El informe post-mortem permite revisar el proyecto con perspectiva una vez finalizado,
identificando aciertos, problemas y decisiones que podrían replantearse en un nuevo ciclo
de desarrollo.

\subsection{Lo que ha ido bien}
\begin{itemize}
    \item La modularización de la memoria y del repositorio ha facilitado mantener alineados el desarrollo técnico y la documentación.
    \item El uso de FDD \cite{palmer2002fdd} ha permitido descomponer el trabajo en funcionalidades manejables y justificar el avance de forma incremental.
    \item La definición temprana de la arquitectura y de la estrategia de control de versiones ha reducido el riesgo de integración entre subsistemas.
\end{itemize}

\subsection{Lo que ha ido mal}
\begin{itemize}
    \item Algunas funcionalidades dependientes de la integración entre sistemas han requerido más iteraciones de las previstas inicialmente.
    \item Parte del contenido documental se ha tenido que reestructurar a medida que evolucionaba el criterio de organización de la memoria.
    \item El equilibrio entre profundidad técnica y nivel de detalle académico ha exigido revisiones sucesivas sobre la forma de presentar el trabajo.
\end{itemize}

\subsection{Discusión}
Si el proyecto se iniciara de nuevo, convendría cerrar antes el esquema final de la
memoria y fijar con mayor precisión la correspondencia entre features, iteraciones y
capítulos. Aun así, el proceso seguido ha permitido obtener una base metodológica y
técnica sólida sobre la que construir la versión final del TFG.

\subsection{Horas reales de la fase de cierre}
La fase de cierre concentró el esfuerzo de consolidación final de la memoria, revisión de
coherencia entre lo implementado y lo documentado, preparación de materiales de entrega y
redacción de conclusiones. El registro horario de esta fase quedó reflejado del siguiente
modo:

\begin{table}[h!]
    \centering
    \begin{coolTable}{lr}{2}
    {Horas de la fase 3}
    \textbf{Horas planificadas de la fase} & 100:00 \\
    \textbf{Horas reales (Adolfo Borrego González)} & 40:26 \\
    \textbf{Horas reales (Germán Sánchez Carmona)} & 49:27 \\
    \textbf{Horas reales totales de la fase} & 89:53 \\
    \textbf{Horas acumuladas planificadas hasta esta fase} & 600:00 \\
    \textbf{Horas acumuladas reales hasta esta fase} & 600:00 \\
    \end{coolTable}
    \caption{Resumen de horas de la Fase 3.}
\end{table}

\section{Trabajos futuros}
Entre las líneas de continuidad del proyecto destacan la ampliación de contenido jugable,
la mejora del sistema de equilibrio y progresión, la profundización en la dualidad entre
personajes y la incorporación de más herramientas de validación y \textit{playtesting}.
Todas estas extensiones podrían desarrollarse como evolución directa del videojuego o como
nuevos trabajos centrados en áreas concretas del sistema.
